\subsection{Setup}
The experiments were conducted using multiple computational setups. 
Training was performed on several Windows-based workstations equipped with NVIDIA A4000 GPUs, using appropriately configured development environments managed with \texttt{conda}. 
In addition, training runs were executed on Google Colab Pro, utilizing NVIDIA A100 and T4 GPUs. 
To ensure availability and persistence of data, results, and trained models, a pipeline integrating GitHub and Google Drive was established. 
For each training run, the two best-performing models as well as the final trained model were stored on Google Drive. 
Experiments were organized using Git branches, and the corresponding notebooks were retrieved from GitHub, modified when necessary, and versioned accordingly.

All training runs and evaluations were logged using the online experiment tracking service \textit{Weights \& Biases (WB)}%
\footnote{Weights \& Biases: Machine learning experiment tracking. https://wandb.ai/site (accessed February 10, 2026).}. 
The platform enabled rapid inspection of intermediate results and provided aggregated metrics to support further optimization.

We implemented several training setups. 
Initially, the goal was to train two models from the residual nets, namely 18-layer and 50-layer ResNet. 
These models served as an initial baseline architectures for subsequent development and comparison. 
During training, hyperparameter sweeps were conducted to identify the optimal learning rate and the corresponding batch size. 
The primary focus was on achieving a high \textit{validation accuracy}. 
The resulting models were retained and reused in later stages of the project.

At this stage, the main objective was to become familiar with the training setup and to establish a functional workflow covering data transfer, model training with validation, and systematic evaluation of results.
